% =============================================================================
%  SECTION 4 — C2: Implicit Indication via Homogenization
%  Slides 4.1–4.8  |  Target time: ~13 minutes
%
%  Corresponds to dissertation Ch. 4 (Contribution 2)
% =============================================================================

\begin{frame}
    \frametitle{C2 Motivation:}
    \framesubtitle{What If the Observation Space Were Invariant?}

    \textbf{A question from C1:}
    \begin{itemize}
        \item In Level-Based Foraging, \(O_i\) incompatibility
            was a significant obstacle.
        \item This raised a natural question: \emph{what if we could make
            the observation space invariant to agent number?}
    \end{itemize}

    \vspace{0.4em}
    \textbf{Surveying the landscape:}
    \begin{itemize}
        \item Parameter sharing unlocks sample efficiency, cross-agent
            generalization, and storage economy — but requires a
            \textbf{common input space}.
        \item Most relevant existing approaches:
            \textbf{one-hot agent IDs}~\cite{terry2020}
                (brittle; no zero-shot transfer) and
            \textbf{separate policies}~\cite{papoudakis2021}
                ($|I|\times$ storage; no cross-type sharing).
        \item Little work addresses invariant representations
            as a direct path to parameter sharing for
            \emph{intrinsically} heterogeneous teams.
    \end{itemize}

\note[note]{Recall difficulty in obs from LBF}

\note[item]{what if we could make it invariant to number of allies?}
\note[item]{Which then naturally leads to invariant to subsets of the observation space}
\note[item]{Further, explain relevant approaches}
\note[note]{separate policies}
\note[note]{one-hot}
\end{frame}


\begin{frame}{C2 Research Questions}
    \begin{itemize}
        \item[\textbf{RQ2.1}] 
        How do input-invariant (homogenized) structures affect learning efficiency?
        \item[\textbf{RQ2.2}] 
        Do input-invariant architectures stabilize performance under
        team-size changes, sensor dropout, and novel agent compositions?
        \item[\textbf{RQ2.3}] 
        What are the computational and implementation costs
        relative to separate per-type policies?
    \end{itemize}
        
\note{About 45 seconds — this is a reference slide.}
\note[item]{Learning Efficiency}
\note[item]{Robustness \& Generalization}
\note[item]{Cost vs.\ Benefit}
\note[note]{Mention the HAPPO baseline explicitly:}
\note[note]{strong baseline for behavioral heterogeneity}
\end{frame}


\subsection{Hypergrid}


\begin{frame}{HyperGrid: Purpose-Built Evaluation Environment}
\begin{columns}
    \begin{column}{0.55\textwidth}
        \textbf{Why a new environment?}
        \begin{itemize}
            \item Existing benchmarks lack \textbf{configurable sensor heterogeneity}.
            \item Need to isolate observation structure as the \emph{independent variable}.
        \end{itemize}
        \vspace{0.4em}

        \textbf{Design:}
        \begin{itemize}
            \item $n$-dimensional discrete gridworld (derived from MiniGrid).
            \item Agent type $=$ sensor channel configuration $C_i \subseteq C$.
            \item Definable dynamics, actions, and rewards for agents.
        \end{itemize}
    \end{column}
    \begin{column}{0.43\textwidth}
        \begin{figure}
            \centering
            \includegraphics[width=0.85\linewidth]{marl-minigrid.png}
            \caption{Two-dimensional rendering of the HyperGrid environment
                    (MiniGrid visual interface).}
            \label{fig:c2-hypergrid}
        \end{figure}
    \end{column}
\end{columns}

\note[note]{Explain why new environement}
\note[note]{more control over parts of the observation space}
\note[note]{kind of as a coding challenge}

\end{frame}

% #TODO: Expand upon Hypergrid implementation

\begin{frame}[shrink=5]{The Implicit Indication Framework}
\vspace{1.5em}
\begin{columns}
    \begin{column}{0.48\textwidth}
        \textbf{Core construction:}
        \begin{itemize}
            \item \textbf{Homogenized observation space:}
                $\widetilde{{O}} = \bigcup_{i \in {I}} {O}_i$,
                union of all agent observation elements.
            \item Each agent populates \emph{only its active channels};
                unused elements are \textbf{zero}.
            \item A \textbf{single shared policy} $\pi_\theta(\cdot \mid \tilde{o})$
                learns to condition on the pattern of populated elements.
        \end{itemize}
    \end{column}
    \begin{column}{0.50\textwidth}
        \begin{figure}
            \centering
            \resizebox{!}{0.75\linewidth}{%
                \input{../Chapters/C2/Figures/rgb_example_open_grid.tex}}
            \caption{An example 2D hypergrid with RGB channels as observation elements.
                The top layer is ground truth, subsequent layers represent observation elements.}
                \label{fig:rgb_hypergrid_example}
        \end{figure}
    \end{column}
\end{columns}

\note[note]{Use "Technicolor Hypergrid" to motivate observation elements}
\note[note]{ground truth is the totality}
\note[note]{but filters represent subsets of access to info}
\end{frame}


\begin{frame}[shrink=5]{The Implicit Indication Framework}
\vspace{1.5em}
\begin{columns}
    \begin{column}{0.48\textwidth}
        \textbf{Required condition:}
        \begin{itemize}
            \item \textbf{Semantic decomposability} — observation elements
                must have type-consistent meanings across agents
                (e.g., ``distance to target'' means the same thing
                for every agent that reports it).
        \end{itemize}
        \vspace{0.3em}

        \textbf{Theoretical grounding:}
        \begin{itemize}
            \item Disjoint injective embeddings preserve
                  representability.
            \item A single function over $\widetilde{\mathcal{O}}$ can
                  replicate any collection of per-type policies.
        \end{itemize}
    \end{column}
    \begin{column}{0.50\textwidth}
        \begin{figure}
            \centering
            % \resizebox{!}{0.75\linewidth}{%
            %     \input{../Chapters/C2/Figures/rgb_example_open_grid.tex}}
            % \caption{An example 2D hypergrid with RGB channels as observation elements.
            %     The top layer is ground truth, subsequent layers represent observation elements.}
            % \includegraphics[width=\linewidth]{disjoint_injective.tex}
            \includegraphics[width=0.75\linewidth]{Figures/disjoint_injective.pdf}
            \label{fig:injective_disjoint}
        \end{figure}
    \end{column}
\end{columns}
\note{    
    \textbf{Beat 2 — The requirement:} ``This only works if the channels mean the same thing
    across agents. We call that semantic decomposability. In practice, it means your sensor
    outputs have to be aligned — distance means distance, heading means heading.''
    \textbf{Beat 3 — The guarantee:} ``Deep Sets tells us that if the embedded regions
    are disjoint, we lose no expressivity. One network can, in principle, do
    everything that $|\mathcal{I}|$ separate networks can do.''
    Point to the channel figure on the right to make it concrete.}
\end{frame}


\begin{frame}[shrink=10]{Training Setup}
    \vspace{2em}
    \begin{columns}
        \begin{column}{0.69\textwidth}
            \textbf{Baseline:} HAPPO~\cite{zhong2024}: separate actor + critic per agent type.
        
            \vspace{0.4em}
            \textbf{Controlled comparison:}
            \begin{itemize}
                \item Identical network architecture and hyperparameters across methods.
                \item Teams of 4 agents, trained for up to 500M steps.
                \item 30 independent seeds per condition.
            \end{itemize}
            \vspace{0.4em}
        
            \textbf{Four team compositions} (varying observation-element coverage):
            \renewcommand{\arraystretch}{1.05}
            \begin{tabular}{ll}
                \emph{Complete:} & $C_i = C$ for all $i$ \\
                \emph{Intersecting-span:} & $\bigcup C_i = C$, overlapping \\
                \emph{Disjoint-span:} & $\bigcup C_i = C$, non-overlapping \\
                \emph{Incomplete:} & $\exists\, c \in C$ not covered by any agent 
            \end{tabular}
        \end{column}
        \begin{column}{0.3\textwidth}
            \begin{figure}
                \scalebox{.65}{
                \subimport{../Chapters/C2/Figures/}{team_compositions.pdf_tex}}
            \end{figure}
        \end{column}
    \end{columns}

\note[note]{We compare to HAPPO}
\note[note]{same parameters as possible}
\note[note]{Happo results in several policies, ours, one}

\note[note]{\textbf{why not explicit indication?}}
\note[note]{Explicit cannot represent the novel configs}

\end{frame}


\begin{frame}{Evaluation Protocol}
    \textbf{Eight evaluation conditions:}
    \begin{columns}
        \begin{column}{0.48\textwidth}
            \begin{enumerate}
                \item Unmodified team (control)
                \item Loss of an agent
                \item Sensor degradation ($-1$ channel)
                \item Sensor improvement ($+1$ channel)
            \end{enumerate}
        \end{column}
        \begin{column}{0.48\textwidth}
            \begin{enumerate}
                \setcounter{enumi}{4}
                \item Reduced team coverage
                \item Increased team coverage
                \item Shuffled channel assignments
                \item Novel spanning set (zero-shot)
            \end{enumerate}
        \end{column}
    \end{columns}

    \vspace{1em}
    \textbf{Selection and metric:}
    \begin{itemize}
        \item Top 15 of 30 trained policies per condition
            (balances noise removal with variance preservation).
        \item Evaluation episodes longer than training episodes
            to further reduce variance.
        \item Primary metric: mean episodic return.
    \end{itemize}


    \note{About 1 minute.
    Frame the eight conditions as a stress test:
    ``We trained under four known configurations.
    Then we asked: what happens when the world changes?
    What if an agent drops out? What if a sensor fails — or improves?
    What if we shuffle who has what? And the hardest test:
    what if we deploy a team composition that was never seen in training?''
    The selection protocol deserves a sentence:
    ``We took the top half of seeds to filter out clearly failed runs,
    then evaluated on longer episodes for tighter confidence intervals.''
    Then transition: ``Now let us see what happened.''}
\end{frame}


\begin{frame}[shrink=15]{Training Results}
\vspace{1em}
\begin{columns}
    \begin{column}{0.45\textwidth}
        \textbf{Comparison:}
        \begin{itemize}
            \item Implicit Indication
            \item HAPPO
            \item Identical network architecture and hyperparameters.
        \end{itemize}

        \vspace{0.5em}
        \textbf{Findings:}
        \begin{itemize}
            \item \textbf{Matched performance}
            \item No statistically significant differences in Training Reward, CPU, RAM,
                or wall-clock time per step.
        \end{itemize}

        \vspace{0.3em}
        \begin{alertblock}{Take-away}
            Same reward, $\frac{1}{|I|}$ the storage.
        \end{alertblock}
    \end{column}
    \begin{column}{0.53\textwidth}
        \begin{figure}
            \centering
            \includegraphics[width=0.8\linewidth]{training_curves.png}
            \caption{Training curves (mean $\pm$ s.d.) for Implicit Indication
                     vs.\ HAPPO across four visibility configurations.
                     30 independent seeds per condition.}
            \label{fig:c2-training}
        \end{figure}
    \end{column}
\end{columns}

    \note{About 2 minutes.
    Not shown on slide is the non-difference in the other measurements}
\end{frame}


\begin{frame}{Robustness: Eight Perturbation Conditions}
    \begin{figure}
        \vspace{1em}
        \begin{columns}
            \begin{column}{0.48\textwidth}
                \centering
                \includegraphics[width=0.9\linewidth]{performance_delta_cropped.png}
            \end{column}
            \begin{column}{0.48\textwidth}
                \begin{alertblock}{Key Result}
                    Policies trained via Implicit Indication are more robust and 
                    exhibit superior handling of novel configurations.
                \end{alertblock}
                \vspace{1em}
                \caption{Performance difference (Implicit Indication $-$ HAPPO)
                        across training configurations (rows) and evaluation
                        conditions (columns). Positive values favor
                        Implicit Indication.}
                \label{fig:c2-delta}
            \end{column}
        \end{columns}
    \end{figure}

    \note{About 2 minutes.
    This is the slide that demonstrates deployment flexibility.
    End by connecting to RQ2.2: ``These results directly answer RQ2.2 —
    yes, the method stabilizes performance under all tested perturbations.''}
\end{frame}

